\documentclass{amsart}

\usepackage[T1]{fontenc}
\usepackage{amssymb}
\usepackage{booktabs}
\usepackage[hidelinks]{hyperref}

\newtheorem{lemma}{Lemma}
\newtheorem{prop}[lemma]{Proposition}
\newcommand{\tablegap}{\par\vspace{0.2cm}}

\title[Analytic verification]{Analytic verification of the uniform
separation estimates}
\author{Jiming Ma}
\author{Fangting Zheng}
\date{July 21, 2026}

\begin{document}

\begin{abstract}
This note supplies the one-variable calculus details used in the uniform
three-layer construction for the right-angled Coxeter groups with
\(p=2,3\) and \(m\geq5\).  It proves that all pole vectors are spacelike,
reduces every non-edge orbit to one representative, and verifies the strict
separation inequality \(H>1\) uniformly in \(m\).  The argument is analytic:
the only numerical enclosures below come from rational Taylor bounds for
trigonometric functions at the two endpoints.
\end{abstract}

\maketitle

\section{Notation and positive quantities}

Fix \(p\in\{2,3\}\), put \(t=2\pi/m\), and allow temporarily any real
\(t\) with
\[
 0<t\leq \frac{2\pi}{5}.
\]
Set
\[
 q=\cos t,\qquad r=\cos\frac tp,
 \qquad U=1-\rho_pt^2,\qquad V=1-\sigma_pt^2,
\]
where
\[
 (\rho_2,\sigma_2)=\left(\frac1{16},\frac1{16}\right),
 \qquad
 (\rho_3,\sigma_3)=\left(\frac{17}{100},\frac{13}{50}\right).
\]
We further write
\[
 E_U=1-qU,\qquad E_V=1-qV,
 \qquad
 P_U=U-1+\frac{E_U}{r},\qquad
 P_V=V-1+\frac{E_V}{r}.
\]
Finally, define
\begin{align*}
 R&=\frac{U(1+q)}{2q},
 &d_B&=R^{-1}=\frac{2q}{(1+q)U},\\
 \Lambda&=\cos\frac t2\sqrt{UV}
 +\frac{\cos(t/(2p))}{r}\sqrt{E_UE_V},
 &\mathcal M&=\sqrt R\,\Lambda
 +\sqrt{R-1}\sqrt{\Lambda^2-1},\\
 d_C&=\mathcal M^{-2},
 &\eta_p&=\cos\frac{2\pi}{p}.
\end{align*}

\begin{lemma}\label{lem:positive}
On \(0<t\leq2\pi/5\), all quantities under square roots in the preceding
definitions are positive.  Moreover,
\[
 0<d_B<1,\qquad \Lambda>1,\qquad \mathcal M>1,
 \qquad P_U>0,\qquad P_V>0.
\]
The Lorentzian squared norms in the construction are
\begin{equation}\label{eq:norms}
 Q_A=\frac{1-q}{q},\qquad Q_B=d_BP_U,\qquad
 Q_C=d_CP_V,\qquad Q_{O_c}=\frac1{\mathcal M^2-1}.
\end{equation}
In particular, all pole vectors are spacelike.
\end{lemma}

\begin{proof}
We have \(0<q<r<1\).  Since
\[
 \frac{13}{50}\left(\frac{2\pi}{5}\right)^2<1,
\]
both \(U\) and \(V\) are positive; hence so are \(E_U\) and \(E_V\).
Also
\[
 \frac{1-q}{1+q}=\tan^2\frac t2>\frac{t^2}{4}>\rho_pt^2,
\]
which is equivalent to \(R>1\) and therefore to \(0<d_B<1\).
Because
\[
 P_U=(1-q)U+\frac{1-r}{r}E_U,
 \qquad
 P_V=(1-q)V+\frac{1-r}{r}E_V,
\]
the two \(P\)-terms are positive.

For later reference, direct differentiation of \(\Lambda\), followed by
multiplication by its positive radical denominator, gives
\begin{equation}\label{eq:lambda-derivative}
 \Lambda_2'(t)>\frac{18}{25}t,
 \qquad
 \Lambda_3'(t)>\frac{6}{25}t.
\end{equation}
To check these two inequalities without decimal arithmetic, use
\(\sin x<x<\tan x\) for the half-angles occurring in the derivative and
insert the rational values of \(\rho_p,\sigma_p\).  After bringing the two
sides to a common positive denominator, the remaining coefficients are
nonnegative and the constant coefficient is positive.  Thus \(\Lambda\) is
strictly increasing.  Its expansion at the origin begins with \(1+O(t^2)\),
so \(\Lambda>1\).

Write \(R^{1/2}=\cosh u\) and \(\Lambda=\cosh v\), where \(u,v>0\).
The addition formula gives \(\mathcal M=\cosh(u+v)>1\).  Substitution in
the four Lorentzian norms yields \eqref{eq:norms}.
\end{proof}

\section{Reduction of the cyclic indices}

For a non-edge pair \(X,Y\), let
\[
 D_{XY}(k)=-\frac{\langle n(X_1),n(Y_{k+1})\rangle}
 {\sqrt{Q_XQ_Y}},
 \qquad H_{XY}(k)=D_{XY}(k)^2.
\]
The sign convention is chosen so that \(D_{XY}(k)>1\) simultaneously
certifies a negative Lorentzian inner product and hyperparallelism.

\begin{lemma}\label{lem:discrete}
For fixed \(p\) and \(t\), it suffices to consider the following cyclic
indices:
\begingroup
\small
\[
 \begin{array}{c@{\qquad}c}
 \toprule
 \text{family}&\text{representative index}\\
 \midrule
 AA,\ AB,\ BC&k=2\\
 AC&k=1\\
 BB,\ CC&k=m\\
 AO_c,\ BO_c&\text{no cyclic index}\\
 \bottomrule
\end{array}
\]
\endgroup
\tablegap
\end{lemma}

\begin{proof}
The denominators are independent of \(k\).  Every unsquared numerator is
a sum of one or two cosine sequences.  The identity
\[
 \cos((k+1)\theta)-\cos(k\theta)
 =-2\sin\frac{(2k+1)\theta}{2}\sin\frac\theta2
\]
shows that its first differences have constant sign on each half of a
cyclic arc, with at most one sign change.  Reflection in the midpoint of an
arc identifies the two halves.  Comparison at the remaining arc endpoints
gives \(k=2\) for \(AA,AB,BC\), \(k=1\) for \(AC\), and \(k=m\) for
\(BB,CC\).  The central pole is fixed by the cyclic symmetry, so the last
two families have no index.
\end{proof}

\section{The eight one-variable functions}

The representatives in Lemma~\ref{lem:discrete} simplify to
\begin{align}
 \Phi_{AA,p}(t)&=(1+2q)^2,\label{eq:phiAA}\\
 \Phi_{AB,p}(t)&=\frac{2(1-q^2)U}{P_U},\label{eq:phiAB}\\
 \Phi_{AC,p}(t)&=
 \frac{q\bigl(\mathcal M-\sqrt{V/q}\bigr)^2}{(1-q)P_V},
 \label{eq:phiAC}\\
 \Phi_{BB,p}(t)&=
 \left(\frac{U-1+\eta_pE_U/r}{P_U}\right)^2,
 \label{eq:phiBB}\\
 \Phi_{BC,p}(t)&=\frac{\Delta_p(t)^2}{P_UP_V},\label{eq:phiBC}\\
 \Phi_{CC,p}(t)&=
 \left(\frac{V-1+\eta_pE_V/r}{P_V}\right)^2,
 \label{eq:phiCC}\\
 \Phi_{AO_c,p}(t)&=\frac{q(\mathcal M^2-1)}{1-q},
 \label{eq:phiAO}\\
 \Phi_{BO_c,p}(t)&=\frac{\Lambda^2-1}{P_U},\label{eq:phiBO}
\end{align}
where
\begin{align}\label{eq:delta}
 \Delta_p(t)
 &=\sqrt{UV}\left(\cos\frac{3t}{2}-\cos\frac t2\right)\notag\\
 &\quad+\frac{\sqrt{E_UE_V}}{r}
 \left(\cos\frac{3t}{2p}-\cos\frac{t}{2p}\right).
\end{align}
Each \(\Phi_{XY,p}\) is the corresponding value of \(H_{XY}\).
Notice that \(\Delta_p(t)<0\).  The identities follow by substituting
\eqref{eq:norms} and using the two \(BC\)-edge equations to subtract the
terms containing \(\cos(t/2)\) and \(\cos(t/(2p))\).  In particular, no
estimate has yet been used in \eqref{eq:phiAA}--\eqref{eq:phiBO}.

For completeness, we give a directly reproducible derivative audit.  The
primitive derivatives are
\begin{align*}
 q'&=-\sin t,&r'&=-\frac1p\sin\frac tp,
 &U'&=-2\rho_pt,&V'&=-2\sigma_pt,\\
 E_U'&=\sin(t)U+2\rho_ptq,
 &E_V'&=\sin(t)V+2\sigma_ptq,\\
 P_U'&=U'+\frac{E_U'}r-\frac{E_Ur'}{r^2},
 &P_V'&=V'+\frac{E_V'}r-\frac{E_Vr'}{r^2}.
\end{align*}
Put
\[
 \alpha=\sqrt{UV},\qquad
 \beta=\frac{\sqrt{E_UE_V}}r.
\]
Then
\begin{align*}
 \frac{\alpha'}\alpha&=\frac12\left(\frac{U'}U+\frac{V'}V\right),\\
 \frac{\beta'}\beta&=\frac12\left(\frac{E_U'}{E_U}
 +\frac{E_V'}{E_V}\right)-\frac{r'}r,\\
 \Lambda'&=\alpha'\cos\frac t2-\frac\alpha2\sin\frac t2
 +\beta'\cos\frac{t}{2p}-\frac\beta{2p}\sin\frac{t}{2p}.
\end{align*}
Moreover,
\begin{align*}
 R'&=\frac{U'}2\left(1+\frac1q\right)-\frac{Uq'}{2q^2},\\
 \mathcal M'&=\frac{R'\Lambda}{2\sqrt R}+\sqrt R\,\Lambda'
 +\frac{R'\sqrt{\Lambda^2-1}}{2\sqrt{R-1}}
 +\frac{\sqrt{R-1}\,\Lambda\Lambda'}{\sqrt{\Lambda^2-1}}.
\end{align*}

Define
\[
 A_0=\frac{1-q}{q},\qquad W=\sqrt{\frac Vq},
 \qquad N_U=U-1+\frac{\eta_pE_U}{r},
 \qquad N_V=V-1+\frac{\eta_pE_V}{r}.
\]
The logarithmic derivatives of the eight positive functions are therefore
\begin{align*}
 \frac{\Phi_{AA,p}'}{\Phi_{AA,p}}
 &=-\frac{4\sin t}{1+2q},\\
 \frac{\Phi_{AB,p}'}{\Phi_{AB,p}}
 &=\frac{2q\sin t}{1-q^2}+\frac{U'}U-\frac{P_U'}{P_U},\\
 \frac{\Phi_{AC,p}'}{\Phi_{AC,p}}
 &=2\frac{\mathcal M'-W'}{\mathcal M-W}
 -\frac{A_0'}{A_0}-\frac{P_V'}{P_V},\\
 \frac{\Phi_{BB,p}'}{\Phi_{BB,p}}
 &=2\left(\frac{N_U'}{N_U}-\frac{P_U'}{P_U}\right),\\
 \frac{\Phi_{BC,p}'}{\Phi_{BC,p}}
 &=2\frac{\Delta_p'}{\Delta_p}-\frac{P_U'}{P_U}-\frac{P_V'}{P_V},\\
 \frac{\Phi_{CC,p}'}{\Phi_{CC,p}}
 &=2\left(\frac{N_V'}{N_V}-\frac{P_V'}{P_V}\right),\\
 \frac{\Phi_{AO_c,p}'}{\Phi_{AO_c,p}}
 &=\frac{q'}q+\frac{q'}{1-q}
 +\frac{2\mathcal M\mathcal M'}{\mathcal M^2-1},\\
 \frac{\Phi_{BO_c,p}'}{\Phi_{BO_c,p}}
 &=\frac{2\Lambda\Lambda'}{\Lambda^2-1}-\frac{P_U'}{P_U}.
\end{align*}
Here \(W'/W=(V'/V-q'/q)/2\), and \(\Delta_p'\) follows immediately
from \eqref{eq:delta}.  These formulas contain all differentiations used
below.

\begin{lemma}[One-variable calculus audit]\label{lem:calculus}
On \(0<t\leq2\pi/5\), the monotonicity patterns are as follows.  The
symbol \(+/-\) means that the derivative changes sign once from positive
to negative, so the unique critical point is a strict maximum.
\begingroup
\small
\[
\begin{array}{c@{\qquad}cccccccc}
\toprule
p&AA&AB&AC&BB&BC&CC&AO_c&BO_c\\
\midrule
2&-&-&-&+/-&-&+/-&+&+\\
3&-&-&-&+&-&+&-&+/-\\
\bottomrule
\end{array}
\]
\endgroup
\tablegap
Consequently, none of the eight functions has an interior minimum.
\end{lemma}

\begin{proof}
Insert the rational pairs \((\rho_p,\sigma_p)\) into the displayed
logarithmic derivatives.  Multiply each expression by the product of its
positive denominators and radical factors.  Use
\[
 \sin t=2\sin\frac t2\cos\frac t2,
 \qquad 1-\cos t=2\sin^2\frac t2,
\]
and the analogous identities with \(t/p\).  The signs in the first three,
the \(BC\), and the \(AO_c\) columns then follow from
\[
 \sin x<x<\tan x\qquad (0<x<\pi/2)
\]
together with \eqref{eq:lambda-derivative}.  The cleared numerators for
\(BB\) and \(CC\) are positive at the left endpoint.  For \(p=2\),
differentiating either cleared numerator once more, and applying the same
half-angle inequalities, gives a strictly negative quantity; its value at
\(2\pi/5\) is negative.  Hence there is exactly one zero and the sign
change is from positive to negative.  For \(p=3\), the two numerators
remain positive up to the right endpoint.  The cleared \(BO_c\) numerator
remains positive for \(p=2\); for \(p=3\), its derivative is strictly
negative and its endpoint signs are opposite, again giving a unique change
from positive to negative.  All multipliers used in clearing denominators
are positive by Lemma~\ref{lem:positive}; hence no sign is reversed.

This calculation is also easy to check directly from the derivative list:
it involves only the four rational constants
\(1/16,17/100,13/50\) and the elementary inequalities just stated.  The
purpose of displaying the logarithmic derivatives is to make the sign
audit independent of any choice of computer algebra system.
\end{proof}

\section{Endpoint certification}

Lemma~\ref{lem:calculus} reduces the lower bounds to \(t=2\pi/5\) and the
limit \(t\to0\).  The endpoint values below are certified by rational
arithmetic.  One may use
\[
 3.141592653589793<\pi<3.141592653589794
\]
and the alternating Taylor bounds
\begin{align*}
 x-\frac{x^3}{3!}
 &<\sin x<
 x-\frac{x^3}{3!}+\frac{x^5}{5!},\\
 1-\frac{x^2}{2!}+\frac{x^4}{4!}-\frac{x^6}{6!}
 &<\cos x<
 1-\frac{x^2}{2!}+\frac{x^4}{4!}
\end{align*}
on the relevant subintervals; additional alternating terms only sharpen
the displayed decimals.  Square-root enclosures are checked by squaring
rational endpoints.

\begin{prop}\label{prop:endpoints}
The minimizing regimes and certified lower bounds are
\begingroup
\small
\[
\begin{array}{c@{\qquad}cc@{\qquad}cc}
\toprule
&\multicolumn{2}{c}{\text{minimizing regime}}
&\multicolumn{2}{c}{\text{certified lower bound}}\\
\text{family}&p=2&p=3&p=2&p=3\\
\midrule
AA&(t=2\pi/5,k=2)&(t=2\pi/5,k=2)&2.618&2.618\\
AB&(t=2\pi/5,k=2)&(t=2\pi/5,k=2)&2.055&2.286\\
AC&(t=2\pi/5,k=1)&(t=2\pi/5,k=1)&1.506&1.0149\\
BB&(t=2\pi/5,k=m)&(t\to0,k=m)&1.559&1.020\\
BC&(t=2\pi/5,k=2)&(t=2\pi/5,k=2)&2.819&2.769\\
CC&(t=2\pi/5,k=m)&(t\to0,k=m)&1.559&1.638\\
AO_c&t\to0&t=2\pi/5&3.375&2.193\\
BO_c&t\to0&t\to0&1.500&1.494\\
\bottomrule
\end{array}
\]
\endgroup
\tablegap
Every entry is strictly larger than \(1\).  The global minimum is attained
for \(p=3,m=5\) in the \(AC\) family, and
\[
 1.0149343<\Phi_{AC,3}(2\pi/5)<1.0149344.
\]
\end{prop}

\begin{proof}
Substitute \(t=2\pi/5\) in
\eqref{eq:phiAA}--\eqref{eq:phiBO} and apply the rational bounds preceding
the proposition.  At the origin, use the second-order expansions of the
same exact formulas.  For example,
\[
 \lim_{t\to0}\Phi_{BB,3}(t)
 =\left(3\rho_3+\frac12\right)^2=1.0201,
\]
and
\[
 \lim_{t\to0}\Phi_{CC,3}(t)
 =\left(3\sigma_3+\frac12\right)^2=1.6384.
\]
For the \(AA\) family,
\[
 \Phi_{AA,p}(2\pi/5)
 =\left(1+2\cos\frac{2\pi}{5}\right)^2
 =\frac{3+\sqrt5}{2}>2.618.
\]
The other substitutions give, before downward rounding,
\begingroup
\small
\[
\begin{array}{c@{\qquad}cc}
\toprule
\text{family}&p=2&p=3\\
\midrule
AB&2.0558126\ldots&2.286\ldots\\
AC&1.5064191\ldots&1.0149343201\ldots\\
BB&1.5597146\ldots&1.0201\\
BC&2.8191746\ldots&2.769\ldots\\
CC&1.5597146\ldots&1.6384\\
AO_c&3.375&2.193\ldots\\
BO_c&1.500&1.494\ldots\\
\bottomrule
\end{array}
\]
\endgroup
\tablegap
The table in the statement records only downward roundings, so the coarse
decimal display is not being used as an equality.  Comparing the two
endpoints in the three \(+/-\) cases selects the minimizing endpoints shown
there.
\end{proof}

\begin{prop}[Uniform separation]
For \(p\in\{2,3\}\) and every integer \(m\geq5\), every non-edge pair in
the uniform three-layer construction has negative Lorentzian inner product
and satisfies \(H>1\).
\end{prop}

\begin{proof}
Lemma~\ref{lem:discrete} reduces each non-edge orbit to one of the eight
functions.  Lemma~\ref{lem:calculus} moves its minimum to an endpoint, and
Proposition~\ref{prop:endpoints} gives a lower bound greater than \(1\).
Before squaring, the representative numerator has the sign encoded in the
definition of \(D_{XY}\); its first-difference proof in
Lemma~\ref{lem:discrete} preserves that sign on the whole orbit.  Hence
\(D_{XY}>1\), proving both assertions.
\end{proof}

\end{document}
